\subsection{Relation to earlier work}
\label{sec:literature}

The correspondence between Markov arithmetic and the geometry of the
punctured torus is classical. Cohn related Markov forms to simple
geodesics~\cite{Cohn1971}, and Reutenauer gave a constructive bijection
between proper Christoffel words and proper Markov
triples~\cite[Theorem~1]{Reutenauer2009}. These correspondences retain a
whole triple or a marked curve. They do not assert that its largest
coordinate, or the trace of the curve alone, determines it. We use the
classical dictionaries to enumerate marked occurrences and retain their
homology throughout the calculation.

The parameter in this paper is attached to a pair of curves. After choosing
one member as a basis element, the other has primitive homology
$(h,hk+t)$ and a prescribed Christoffel gap pattern. This differs from
fixing a numerator in one global Farey parametrization. The ordering
theorems of Lagisquet, Pelantov{\'a}, Tavenas and Vuillon~\cite{LagisquetEtAl2021}
and of Lee, Li, Rabideau and Schiffler~\cite{LeeEtAl2023} compare Markov
values along specified directions in that global parametrization.
Our bounds allow the integral character of the chosen basis to vary.
Changing the basis transforms both curves and their character; it cannot
be used to reset the first curve while leaving its trace equation unchanged.

Effective finiteness on character varieties is also established. Whang
proves effective bounds for integral points on geometrically irreducible
nonintegrable curves in relative character varieties, and describes the
multitwist exceptions in the more general nondegenerate
case~\cite[Theorems~1.2--1.3]{Whang2020}. Applied to a fixed trace-equality
curve, this method requires verification of the exceptional hypotheses
component by component. The present argument gives explicit bounds in
the pair parameter $h$ and a direct finite comparison domain, without
computing those algebraic curves. Its quantitative conclusion is
$O(h^5)$ trace comparisons on $O(h^3)$-bit integers. This is a bound in
the numerical value of $h$, and does not bound $h$ itself.

The geometric part uses equally classical structure. The length of a
primitive simple class on a punctured torus is determined by its primitive
homology vector, and these lengths extend to a norm~\cite{McShaneRivin1995}.
McShane and Parlier prove that, with the boundary length fixed, the locus
where two distinct simple curves have equal length is an embedded
path~\cite[Theorem~1.3]{McShaneParlier2008}. We use this theorem
for the reflection cases. The remaining arithmetic classifications retain
the full marking and test the actual finite isometry group of the modular
torus. Homological equivalence under an arbitrary change of basis is not
an isometry test at a fixed metric.

The comparison model for the analytic estimates has a useful interpretation
on the Cayley cubic. For full traces, put
$$
F(X,Y,Z)=X^2+Y^2+Z^2-XYZ.
$$
The Fricke identity gives $F=\operatorname{tr}[A,B]+2$, so the cusped
Markov character lies on $F=0$~\cite{Goldman2003}. Commuting diagonal
matrices instead give
$$
(X,Y,Z)=(u+u^{-1},v+v^{-1},uv+(uv)^{-1}),\qquad F=4.
$$
The torus cover and monomial action on this second level are
classical~\cite[Sections~1.5 and~2.1]{CantatLoray2009}. Taking
$(u,v)=(s^a,s^b)$ turns a Vieta move into subtraction of the exponents.
The same polynomial move preserves the actual Markov level separately.
Our estimates control the discrepancy along these two evolutions;
after division of the traces by three, their level difference is $4/9$.
There is no identification of the two character surfaces in this argument.

Positive trace expansions and quadratic norm gaps have substantial prior
uses. Fisac's necklace formula already records the dependence of a trace
on cyclic runs~\cite[Theorem~1.1 and Section~4]{Fisac2025}.
Altassan and Luca use Binet expansions and quadratic norms to bound
solutions of Markov-type equations in Lucas
sequences~\cite[Theorem~1.1 and Lemmas~2.1--2.2]{AltassanLuca2021}.
Their sequences have specified initial values. Here a centered Markov ray
has coefficients satisfying $cd=p^2/(9p^2-4)$, and these coefficients
vary with the integral character. The needed estimates therefore retain
both coefficients, all mixed matrix paths and the actual Vieta mutations.
At minimum index, mechanical balance controls the mixed paths, while
spectral-angle errors propagate through the Euclidean algorithm with
polynomial amplification. The contradiction is an actual positive
integer forced below one.

Finally, arithmetic restrictions on the largest Markov entry, such as
Button's ideal-theoretic criteria~\cite{Button2001}, organize partial
uniqueness by a different parameter. The applications here concern complete
families of relative curve positions: occurrence counts two through six,
and all intersections at most twenty-four. No upper bound on the common
Markov value is assumed in these classifications. These classifications and
the effective decision procedure leave open whether a nonisometric equality
occurs at some larger occurrence count.
